% Chapter 4 (results half) -- Detection Results, Ablations, Range, Runtime
% Sections 4.2--4.5. Standalone, compilable draft (FPT-template article class +
% booktabs), sibling to sec_4_1_evaluation_protocol.tex.
%
% MERGE/NUMBERING NOTE: this file opens with \section{Detection Evaluation} for
% standalone compilation and numbering, but carries a DISTINCT label
% (\label{ch:detection-eval-results}) so it does not clash with sec_4_1's
% \label{ch:detection-eval}. At final assembly the assembler drops this
% \section line and the subsections below fold under sec_4_1's chapter, so that
% 4.1 (protocol) ... 4.5 (runtime) are one chapter. Forward/back references to
% sibling sections use typed numerals ("Section~\ref{sec:eval-protocol}", "Chapter~\ref{ch:recall-bottleneck}") exactly as
% sec_4_1 does, because the two files do not share a label namespace when built
% separately.
%
% Honesty-phrasing rules applied (personal_agenda P1):
%   - "mean IoU" always qualified as "of matched detections"
%   - recall denominator restated / cross-referenced (>=10 surviving points)
%   - seq 08 = classifier model-selection set, not a held-out test set
%   - negative results reported AS FINDINGS, not apologised for
%     (C5 synthetic-pretraining redundant; clustering "adopt nothing")
%   - the distance table is per-frame / track-filter-OFF: stated as a deviation,
%     read as a trend, NOT a restatement of the 0.730 headline
%   - every number traces to a finding / output artifact (cited Finding~\#NN)
%
\section{Detection results}
\label{sec:results-detection}

The headline operating point is the sequence~08 row of Table~\ref{tab:headline}
(Section~\ref{sec:eval-protocol}): precision \num{0.905}, recall \num{0.730}, $F_1$ \num{0.808}, and
a mean IoU of matched detections of \num{0.912}, over
\num{25478}~/~\num{2676}~/~\num{9444} TP~/~FP~/~FN. Two properties of this
operating point organise the rest of the chapter, and both are visible directly
in the per-frame record rather than only in the pooled totals.

First, the pipeline is \emph{precision-saturated}. Across the full sequence the
false-positive count is low and stable --- about \num{0.7} per frame
(\num{2676} over \num{4071} frames), independent of how many cars are in view.
Figure~\ref{fig:bev-detections} shows this spatially --- detections track the
true parked and moving cars along the street, with few spurious boxes --- and
Figure~\ref{fig:timeseries} shows it over time: the FP trace is low and
stationary while the number of ground-truth cars per frame swings widely.

Second, the pipeline is \emph{recall-limited}, and the limit is structural. The
recall shortfall traces to a measured clustering failure: HDBSCAN splits large,
close cars into fragments that no longer match a single instance
(Finding~\#23), and the split rate is highest exactly where cars are densest and
nearest (\SI{66}{\percent} at \SIrange{0}{10}{\metre}, Finding~\#24).
Section~\ref{sec:results-distance} exhibits the resulting near-range recall dip
directly, and the time series (Figure~\ref{fig:timeseries}) shows per-frame
recall tending to fall in the denser frames while false positives stay flat. We
root-cause the splitting mechanism in Chapter~\ref{ch:recall-bottleneck}; here we establish only that
recall, not precision, is the dominant error mode, and that it does not come from
the classifier misfiring (false positives stay flat regardless).

\begin{figure}[htbp]
  \centering
  \includegraphics[width=\linewidth]{seq08_bev_detections.png}
  \caption{Bird's-eye view of detections on a representative sequence~08 frame.
    Detected car clusters (coloured) are overlaid on the ground-removed point
    cloud; ego is at the origin. Detections concentrate on true vehicles, with
    few spurious clusters --- the precision-saturated behaviour reported in
    Table~\ref{tab:headline}.}
  \label{fig:bev-detections}
\end{figure}

\begin{figure}[htbp]
  \centering
  \includegraphics[width=\linewidth]{seq08_timeseries.png}
  \caption{Per-frame detection behaviour over sequence~08. False positives (lower
    trace) stay low and flat --- under one per frame --- regardless of scene
    content, while per-frame recall tends to fall in frames with more
    ground-truth cars. The structural recall limit is root-caused in Chapter~\ref{ch:recall-bottleneck}
    (HDBSCAN splitting, Findings~\#23/\#24).}
  \label{fig:timeseries}
\end{figure}

\section{Ablations}
\label{sec:results-ablations}

This section decomposes the headline operating point into the contribution of
each learned or post-hoc stage. Every ablation row is produced under the frozen
protocol of Section~\ref{sec:eval-protocol} unless a deviation is stated explicitly.

\subsection{Stage decomposition: classifier and track-filter contributions}
\label{sec:ablation-stage}

Table~\ref{tab:ablation-stage} isolates the two filtering stages that sit on top
of geometric clustering. The geometric-only row uses clustering plus the
geometric size/shape filters with \emph{no} learned classifier and \emph{no}
track filter (Finding~\#47, task~B6, like-for-like under the promoted
configuration); the full row is the headline.

\begin{table}[H]
  \centering
  \caption{Stage ablation on the full sequence~08 (frozen configuration). Rows
    add one filtering stage at a time on top of geometric clustering. The
    per-frame classifier is the precision mechanism (it removes $\sim$\num{149000}
    junk clusters) but is conservative on borderline cars (recall drops); the
    offline track filter then recovers recall by temporal voting while raising
    precision further. Mean IoU is over matched detections only. Geometric-only
    and the middle row are per-frame (no track filter); the full row uses the
    offline track filter.}
  \label{tab:ablation-stage}
  \setlength{\tabcolsep}{5pt}
  \begin{tabular}{lccccc}
    \toprule
    Configuration & Precision & Recall & $F_1$ & Mean IoU & TP / FP / FN \\
    \midrule
    Geometric only                & 0.149 & 0.775 & 0.250 & 0.907 & \num{27051} / \num{154612} / \num{7871} \\
    \;+\,classifier              & 0.820 & 0.677 & 0.741 & 0.921 & \num{23629} / \num{5204} / \num{11293} \\
    \;+\,classifier, track filter & 0.905 & 0.730 & 0.808 & 0.912 & \num{25478} / \num{2676} / \num{9444} \\
    \bottomrule
  \end{tabular}
\end{table}

The decomposition tells a two-step story (Finding~\#47). The \textbf{classifier
is the precision mechanism}: geometric clustering under the promoted coarse-voxel
setting emits many mostly-junk clusters (precision \num{0.149}), and the per-frame
classifier removes about \num{149000} of them across the sequence
(FP \num{154612}~$\rightarrow$~\num{5204}), lifting precision from \num{0.149} to
\num{0.820}. It pays for this with recall --- being conservative, it discards
borderline true clusters too, so recall falls from \num{0.775} to \num{0.677}.
The \textbf{track filter then recovers that recall while improving precision
again}: linking detections across frames and accepting only consistent tracks
lifts precision to \num{0.905} \emph{and} recall to \num{0.730} (TP
\num{23629}~$\rightarrow$~\num{25478}), because a car seen across many frames is
kept even where a single frame classified it as unknown, while transient
false positives fail the track-length test. The net of both stages is still a
recall trade --- geometric-only recall (\num{0.775}) exceeds the full-pipeline
recall (\num{0.730}), roughly \num{1600} true positives forgone for the
$\sim$\num{152000}-false-positive removal --- but that trade runs through a
recall dip and a partial recovery, not a simple monotone loss. We report it
rather than hiding it: the system is deliberately tuned for precision, and the
recall it forgoes here is small next to the recall it never had (Chapter~\ref{ch:recall-bottleneck}). The
pre-promotion geometric-only figure ($F_1$~\num{0.305}, Finding~\#18) is
superseded by this like-for-like \num{0.250} and kept only as a historical
record.

\subsection{Redundancy of synthetic pretraining (Stage~A ablation)}
\label{sec:ablation-stagea}

The classifier ships as a binary car/not-car network trained \emph{from scratch}
on mined real clusters (Stage~B). The project proposal had assumed a two-stage
recipe in which a synthetic-ShapeNet Stage~A pretraining seeds Stage~B. We tested
whether that prior is necessary and found it is not --- a negative result we
report as a finding, because it changes the recommended recipe.

At the classifier level, on the sequence~08 validation clusters, the from-scratch
network matches the pretrained one: macro $F_1$ \num{0.9285} from scratch versus
\num{0.9225} pretrained (Finding~\#25; the two are within run-to-run noise, with
scratch marginally ahead). At the pipeline level a small precision edge for the
pretrained checkpoint appeared on a 100-frame sample (FP \num{27} versus
\num{62}), but it was uncontrolled (different training seed and code) and
\textbf{did not survive at full scale}: on the full sequence~08 the two
checkpoints are identical in $F_1$ to three decimals (Finding~\#31). We therefore
drop Stage~A from the production pipeline and keep it as ablation material.

The reason the synthetic prior is redundant --- and cannot simply be added for
safety --- is a total, symmetric domain gap between synthetic and real point
clusters (Finding~\#30). Table~\ref{tab:cross-domain} gives the cross-domain
transfer matrix: a model trained on one domain and evaluated on the other scores
a car-class $F_1$ of \num{0.000} in \emph{both} off-diagonal directions, and
fine-tuning on real data erases whatever the synthetic pretraining had learned.
Given \num{420000} real training clusters, the real data does essentially all the
work; the synthetic prior neither helps nor, once fine-tuned, survives.

\begin{table}[H]
  \centering
  \caption{Cross-domain classifier transfer (car-class $F_1$; Finding~\#30). Each
    model is trained on one domain and evaluated on the other's validation set;
    diagonal cells (italic) are in-domain, for scale. Both off-diagonal cells
    collapse to \num{0.000} --- against strong in-domain performance
    (\num{0.999}, \num{0.885}) --- so synthetic and real clusters are mutually
    out-of-distribution and synthetic pretraining transfers nothing to real
    inference.}
  \label{tab:cross-domain}
  \begin{tabular}{lcc}
    \toprule
    & \multicolumn{2}{c}{Evaluated on} \\
    \cmidrule(lr){2-3}
    Trained on & Synthetic & Real \\
    \midrule
    Synthetic (Stage~A) & \emph{0.999} & 0.000 \\
    Real (Stage~B)      & 0.000        & \emph{0.885} \\
    \bottomrule
  \end{tabular}
\end{table}

\subsection{Clustering method: cost and benefit of HDBSCAN}
\label{sec:ablation-clustering}

The clustering stage is HDBSCAN. Because it is by far the most expensive stage
(Section~\ref{sec:results-runtime}), we benchmarked it against two cheaper
density/connectivity clusterers --- DBSCAN and Euclidean-tolerance clustering ---
at the same voxel resolution and under an identical pipeline
(Table~\ref{tab:clustering}, Finding~\#43). \textbf{Deviation:} this benchmark is
per-frame with the track filter \emph{off}, because it strides across the
sequence and the tracker requires contiguous frames; the numbers are therefore
comparable to each other but not to the track-filtered headline.

\begin{table}[H]
  \centering
  \caption{Clustering-method benchmark at fixed voxel resolution
    ($\SI{0.10}{\metre}$), identical pipeline, per-frame / track-filter-off
    (Finding~\#43). HDBSCAN wins $F_1$ on both sequences, entirely through
    recall; precision is essentially method-independent. The alternatives are
    5--7$\times$ faster. Mean IoU is over matched detections; the timing column is
    the mean per-frame clustering time.}
  \label{tab:clustering}
  \begin{tabular}{llccccc}
    \toprule
    Seq & Method & Precision & Recall & $F_1$ & Mean IoU & Clustering (mean ms) \\
    \midrule
    \multirow{3}{*}{00}
      & HDBSCAN   & 0.962 & \textbf{0.731} & \textbf{0.831} & 0.968 & 391 \\
      & DBSCAN    & 0.962 & 0.637 & 0.767 & 0.964 & 59 \\
      & Euclidean & 0.967 & 0.636 & 0.768 & 0.982 & 78 \\
    \midrule
    \multirow{3}{*}{08}
      & HDBSCAN   & 0.813 & \textbf{0.671} & \textbf{0.735} & 0.923 & 451 \\
      & DBSCAN    & 0.848 & 0.565 & 0.678 & 0.927 & 69 \\
      & Euclidean & 0.838 & 0.576 & 0.683 & 0.950 & 93 \\
    \bottomrule
  \end{tabular}
\end{table}

The gap is entirely a recall gap --- precision is within \num{0.035} across all
methods on either sequence, while HDBSCAN's recall exceeds the alternatives by
\numrange{9}{11} points. An $\varepsilon$-sweep confirms no fixed radius closes
it: the best fixed-radius setting (Euclidean at $\varepsilon=\SI{0.4}{\metre}$)
reaches $F_1$ \num{0.800} on sequence~00, still short of HDBSCAN's \num{0.831}.
HDBSCAN's variable-density criterion recovers cars that a single radius cannot,
which is exactly the recall margin that matters here. We therefore
\textbf{adopt nothing} from the benchmark and keep HDBSCAN in production, paying
its runtime cost knowingly; DBSCAN and Euclidean remain available as
command-line options for reproduction.

\section{Detection range}
\label{sec:results-distance}

Recall is not uniform in range, and the way it varies is not the naive
``far~=~worse'' monotone one might expect. Table~\ref{tab:distance} stratifies
detection by the range of each cluster's centroid from the ego sensor, on a
stride-20 sample of sequence~08 under the frozen configuration.
\textbf{Deviation:} like the clustering benchmark, this table is
\emph{per-frame with the track filter off} (striding breaks the tracker), so its
pooled recall (\num{0.671}) is a per-frame figure and does \emph{not} restate the
track-filtered headline \num{0.730}; read the table as a range \emph{trend}, not
as a re-derivation of the headline.

\begin{table}[H]
  \centering
  \caption{Detection stratified by range-to-ego on sequence~08 (stride-20 sample,
    frozen configuration, per-frame / track-filter-off). Recall is
    \emph{non-monotonic}: it dips in the nearest bin, peaks at
    \SIrange{10}{20}{\metre}, then falls with range. The nearest-bin dip is
    over-segmentation (Finding~\#23), the far decline is sparsity. Mean IoU is
    over matched detections. Summing the two near bins recovers the
    \SIrange{0}{20}{\metre} figure ($R=\num{0.829}$) comparable to the earlier
    Finding~\#5.}
  \label{tab:distance}
  \begin{tabular}{lccccc}
    \toprule
    Range (m) & Precision & Recall & $F_1$ & Mean IoU (matched) & TP / FP / FN \\
    \midrule
    0--10   & 0.896 & 0.787 & 0.838 & 0.902 & \num{311} / \num{36}  / \num{84}  \\
    10--20  & 0.855 & \textbf{0.862} & 0.859 & 0.926 & \num{432} / \num{73}  / \num{69}  \\
    20--40  & 0.762 & 0.549 & 0.638 & 0.934 & \num{400} / \num{125} / \num{329} \\
    40+     & 0.393 & 0.196 & 0.262 & 0.943 & \num{22}  / \num{34}  / \num{90}  \\
    \midrule
    Pooled  & 0.813 & 0.671 & 0.735 & 0.923 & \num{1165} / \num{268} / \num{572} \\
    \bottomrule
  \end{tabular}
\end{table}

The counter-intuitive result is the \textbf{nearest-bin recall dip}: recall at
\SIrange{0}{10}{\metre} (\num{0.787}) is \emph{lower} than at
\SIrange{10}{20}{\metre} (\num{0.862}), even though near cars carry far more
points. This is not sparsity --- it is over-segmentation. Finding~\#24 measures
that HDBSCAN splits the majority of the very-near cars (a split rate of
\SI{66}{\percent} at \SIrange{0}{10}{\metre}, falling to \SI{4}{\percent} at
\SIrange{30}{50}{\metre}): a large, close car dissolves into several clusters,
none of which matches the whole instance, so it is counted as a miss. The nearest bin
also has the \emph{lowest} mean IoU of matched detections (\num{0.902}),
consistent with fragments eroding the overlap even on the cars that do match. We
develop this splitting mechanism, and why post-hoc fragment-merging fails to
recover it, in Chapter~\ref{ch:recall-bottleneck}.

Beyond the peak the decline is ordinary sparsity: recall falls from \num{0.862}
at \SIrange{10}{20}{\metre} to \num{0.549} at \SIrange{20}{40}{\metre} and
\num{0.196} past \SI{40}{\metre}, as cars thin to too few surviving points to
cluster or to clear the \num{10}-point eligibility rule (Section~\ref{sec:eval-protocol}).
Interestingly, mean IoU of matched detections \emph{rises} with range
(\num{0.902}~$\rightarrow$~\num{0.943}): the far cars that are detected at all are
compact, single-cluster instances, so their overlap is clean even as their number
collapses. The earlier precision-only, seq-00, pre-promotion range analysis
(Finding~\#5) is superseded by this on-frozen-config sequence-08 table and is
retained only as a historical note.

\section{Runtime}
\label{sec:results-runtime}

Table~\ref{tab:runtime} breaks the per-frame cost down by stage, measured over
the full sequence~08 under the shipped (promoted) configuration
(\num{4068} timed frames, RTX~3070~Ti / Ryzen~7~7800X3D; Finding~\#34;
per-stage figures from \texttt{timing\_seq08\_full\_n4068\_combined.json}). The breakdown shows where
the time goes: the two classical geometric stages --- ground-plane RANSAC and,
above all, HDBSCAN clustering --- dominate, together taking roughly
\SI{68}{\percent} of the frame; adding preprocessing and the geometric filter,
the classical CPU stages account for about \SI{94}{\percent}, leaving the single
learned component --- the classifier --- at only $\sim$\SI{5.4}{\percent} (the
centroid/Hungarian tracker is a classical assignment step, not learned). The
classifier is \SI{33.7}{\milli\second} for the $\sim$\num{45} clusters in a frame
($\sim$\SI{0.75}{\milli\second} each) and the tracker is negligible.

\begin{table}[H]
  \centering
  \caption{Per-frame runtime by stage on the full sequence~08 under the shipped
    (promoted) configuration (\num{4068} timed frames; mean ms; Finding~\#34).
    Clustering alone is nearly \SI{60}{\percent} of the frame. The total is
    $\approx\SI{623.5}{\milli\second}$ ($\approx$\num{1.6}~fps) on this full run;
    a stride-20 sample of the same configuration estimates
    $\approx\SI{650.6}{\milli\second}$ (the two agree to within the sampling
    variation of the expensive stages). The pre-optimisation configuration
    totalled \SI{921.3}{\milli\second}. Percentages are of the
    \SI{623.5}{\milli\second} total.}
  \label{tab:runtime}
  \begin{tabular}{lcc}
    \toprule
    Stage & Mean (ms) & Share \\
    \midrule
    Data load                       & 0.7   & 0.1\% \\
    Preprocessing (z-filter, denoise, voxel) & 125.9 & 20.2\% \\
    Ground removal (RANSAC)         & 56.2  & 9.0\% \\
    Clustering (HDBSCAN)            & 370.4 & 59.4\% \\
    Geometric filter                & 36.2  & 5.8\% \\
    Classification                  & 33.7  & 5.4\% \\
    Tracking                        & 0.3   & 0.1\% \\
    \midrule
    Total                           & 623.5 & 100\% ($\approx$1.6 fps) \\
    \bottomrule
  \end{tabular}
\end{table}

The promoted configuration is itself the product of a runtime optimisation that
targeted the dominant classical stages directly (Finding~\#34): reducing the
RANSAC iteration count and coarsening the clustering voxel to \SI{0.10}{\metre}
cut the per-frame total by roughly a third, from a pre-optimisation
\SI{921.3}{\milli\second} to \SI{623.5}{\milli\second} on the full sequence
(\SI{650.6}{\milli\second} on the stride-20 sample), and the coarser voxel also
\emph{improved} detection --- it is the promoted configuration used throughout
this chapter. A sub-\SI{400}{\milli\second} target was shown to be unreachable
through the authorised optimisation tiers --- the post-ground object cloud is
intrinsically sparse, so coarser voxelisation compresses it little and the Python
HDBSCAN retains a floor of $\sim$\SI{185}{\milli\second} --- and the achieved
operating point was accepted, real-time operation being out of scope (user
decision, \num{2026}-07-23). Point-completion, when it runs, adds
$\sim$\SI{13}{\milli\second} per completed car on top of detection and is not part
of the detection budget above.
